\clearpage
\fancyhead[L]{\small INSTRUMENTE.}
\hypertarget{native-429-part-lower}{}
\begin{center}{\Large INSTRUMENTE.}\end{center}
\begin{center}
Astronomisches Jahrbuch für 1822. Seite 235..241. Berlin 1819.
\end{center}

\par\indent Den Cometen habe ich zuerst am 2ten Juli in Lauenburg beobachtet, an welchem Abend derselbe daselbst von dem Lieutenant \textsc{Zartmann} aufgefunden wurde. Ich habe indessen diese Beobachtungen (durch Distanzen von Sternen mit Spiegelsextanten gemessen) nicht scharf reducirt, da ich bald nachher erfuhr, dass der Comet um dieselbe Zeit auch schon an andern Orten entdeckt und im Meridian beobachtet war. Nach meiner Zurückkunft habe ich ihn hier bei seiner untern Culmination im Meridian beobachtet; zuerst bis zum 26sten Juli am \textsc{Repsold}schen Kreise, wo jedoch die Antritte an die feinen Fäden anfingen, sich schwer beobachten zu lassen, während die Zenithdistanzen noch mehr Schärfe zu erlauben. 
Da ich inmittelst das \textsc{Reichenbach}sche Mittagsfernrohr (welches bei beträchtlich grösserer optischer Kraft etwas stärkere Fäden hat) in beobachtungsfertigen Stand gesetzt hatte, so beobachtete ich die geraden Aufsteigungen vom 27sten Juli an mit diesem Instrument; während Hr. Prof. \textsc{Harding} noch die Zenithdistanz am \textsc{Repsold}schen Kreise einstellte. Seit dem 4ten August habe ich noch verschiedene Abende (zuletzt am 21sten August) den Cometen am Mittagsfernrohr gesehen; da aber die \SourceLeaf{430}
Antritte mehr geahnt als wirklich beobachtet waren, so unterdrücke ich sie lieber, als keines Vertrauens würdig. Inzwischen bestätigen sie doch die fortdauernde gute Uebereinstimmung der \textsc{Dircksen}schen Elemente.

\Needspace{18\baselineskip}
\begin{center}
\emph{Beobachtungen des Cometen.}
\end{center}

\begin{center}
\begin{tabular}{l|c|c|c}
1819 & Mittlere Zeit & & \\
\hline
Juli 19 & $11^{\mathrm u}\,56^{\mathrm m}\,53^{\mathrm s}$ & $116^\circ\,7'\,26''$ & $+51^\circ\,52'\,12''$ \\
21 & $11\,53\,33$ & $117\,15\,34$ & $51\,54\,58$ \\
22 & $11\,51\,53$ & $117\,49\,34$ & $51\,54\,56$ \\
26 & $11\,44\,14$ & $119\,51\,10$ & $51\,48\,49$ \\
27 & $11\,42\,8$ & $120\,18\,43$ & --- --- --- \\
28 & $11\,40\,2$ & --- --- --- & $51\,43\,46$ \\
30 & $11\,35\,40$ & $121\,38\,52$ & $51\,37\,23$ \\
Aug. 3 & $11\,26\,24$ & $123\,16\,9$ & $51\,23\,25$ \\
4 & $11\,23\,59$ & $123\,38\,55$ & $51\,19\,46$ \\
\end{tabular}
\end{center}

\Needspace{13\baselineskip}
\par\indent Herr \textsc{Dircksen} hat nach diesen und den ersten Beobachtungen des Herrn Prof. \textsc{Struve} in Dorpat folgende parabolische Elemente berechnet.

\begin{center}
\begin{tabular}{ll}
Zeit der Sonnennähe & 1819 Juni 28.\ $17^{\mathrm n}\,49^{\mathrm m}\,9^{\mathrm s}$ Göttinger Zeit \\
Länge der Sonnennähe & . . . . . . . . . . .\ $287^\circ\,6'\,24''8$ \\
Logarithm des Abstandes & . . . . . . . . . .\ $9.5330800$ \\
Länge des Knoten & . . . . . . . . . . . .\ $273^\circ\,42'\,8''8$ \\
Neigung der Bahn & . . . . . . . . . . . .\ $80\,45\,12.0$ \\
\multicolumn{2}{l}{Bewegung rechtläufig.}
\end{tabular}
\end{center}

\Needspace{17\baselineskip}
\par\indent Die Vergleichung obiger Beobachtungen mit diesen Elementen giebt folgende Unterschiede:

\begin{center}
\begin{tabular}{l|c|c}
 &  Ger. Aufst.  &  Abweich.  \\
\hline
Juli 19  &  $-15''5$  &  $-3''5$  \\
21  &  $+39.7$  &  $+5.2$  \\
22  &  $-30.0$  &  $+9.5$  \\
26  &  $-37.9$  &  $+26.5$  \\
27  &  $-9.8$  &  --- ---  \\
\hline
  &  Ger. Aufst.  &  Abweich.  \\
\hline
 Juli 28  &  --- ---  &  $+4''0$  \\
 30  &  $-0''2$  &  $-19.2$  \\
 Aug. 3  &  $-9.1$  &  $-0.3$  \\
 4  &  $+3.8$  &  $-2.9$  \\
\end{tabular}
\end{center}

\Needspace{18\baselineskip}
\par\indent Die Vergleichung einiger früheren auswärtigen Beobachtungen giebt folgende Unterschiede:

\begin{center}
\begin{tabular}{l|c|c|l}
Juli 2  &  $-25''1$  &  $-16''6$  &  Bode  \\
3  &  $-1.0$  &  $-1.6$  &  Struve  \\
3  &  $-1.8$  &  $+44.4$  &  Pond  \\
4  &  $-17.3$  &  $+19.6$  &  Struve  \\
4  &  $-48.3$  &  $-65.8$  &  Bode  \\
5  &  $+12.1$  &  $+18.8$  &  Struve  \\
\hline
 Juli 5  &  $-19''1$  &  $-86''7$  &  Pond  \\
 6  &  $-16.6$  &  $+18.0$  &  Struve  \\
 6  &  $-20.8$  &  $-19.9$  &  Bode  \\
 7  &  $-18.1$  &  $+17.3$  &  Struve  \\
 11  &  $+41.1$  &  $-24.4$  &  Bode  \\
 12  &  $-96.3$  &  $-6.3$  &  Bode  \\
\end{tabular}
\end{center}

\par\indent Das sechsfüssige Mittagsfernrohr von \textsc{Reichenbach} ist nun seit Ende Juli in täglichem Gebrauche; eine kurze Nachricht von diesem Instrumente wird Ihnen nicht unlieb sein. Das Objectiv von \textsc{Frauenhofer}
 hat 52 Pariser Linien Oeffnungen; 4 Oculare vergrössern (nach meiner erst vorläufigen Schätzung) 80, 100, 140, 200 mal. Ich bediene mich gewöhnlich der stärksten, wobei freilich eine sorgfältige Einstellung nöthig ist, um die Sterne nicht zu verfehlen, da der Halbmesser des Gesichtsfeldes nur $4\tfrac{1}{2}$ Minute hält. Vertikalfäden sind 7 eingezogen, deren Intervalle von Aequatorsternen binnen $10^{\mathrm s}$ durchlaufen werden. An diese kleinen Intervalle gewöhnt man sich, nach einiger Uebung, leicht; es geht ganz füglich an, während der Zwischenzeit die Beobachtungen aufzuschreiben, das Ocular zu schieben, \SourceLeaf{431}
und früh genug fertig zu sein, um mit völliger Ruhe die zwei Schläge der Uhr zu erwarten, zwischen
welche der nächste Antritt fällt. Nur muss man gewöhnt sein, während dieser Operationen die Schläge
der Uhr immer sicher fortzuzählen, falls man sie nicht durch einen Gehülfen zählen lassen kann. Sobald
man diese Fertigkeit erworben hat, findet man es überaus angenehm, in kurzer Zeit eine grössere
Menge von Bestimmungen zu erhalten. Den Beweis, dass unter jener Voraussetzung die grössere Anzahl und Nähe der Fäden der Genauigkeit der einzelnen Beobachtungen keinen Eintrag thun, liefern
meine Erfahrungen bereits selbst. Ich finde aus fast 300 beobachteten Antritten von $\gamma$, $\alpha$ und $\beta$ im
Adler den wahrscheinlichen Fehler \emph{eines} Antritts $= 0^{\mathrm{s}}090$, noch etwas kleiner, als bei meinen Beobachtungen am \textsc{Repsold}schen Kreise, wo die 5 Fäden über $16^{\mathrm{s}}$ in Zeit von einander abstehen. — Die Stellung des Fernrohrs geschieht vermittelst eines kleinen an der Seite des Rohrs, nahe beim Ocular, befestigten auf Viertelsgrade getheilten Kreises, auf dem sich eine Alhidade mit einer Libelle bewegt.
Der Vernier an dieser Alhidade giebt einzelne Minuten an. Diese Einrichtung ist der am neuen Greenwicher Mittagsfernrohr ähnlich; sie gewährt den grossen Vortheil, dass man bei umgelegtem Instrument die Stellung mit gleicher Leichtigkeit machen kann, auch bringt man das Rohr, wenn die Alhidade einmal gestellt ist, schneller wieder in die rechte Lage, falls es durch einen Zufall etwas verrückt
wird. Dagegen scheint die Stellung, da sie \emph{zwei} Operationen nöthig macht, ein wenig mehr Zeit zu
erfordern, als bei der gewöhnlichen Einrichtung; auch kann man bei Nacht nicht gut ohne einen Gehülfen fertig werden. Hr. \textsc{von Reichenbach} wird mir jedoch noch einen zweiten Kreis mit etwas veränderter Einrichtung liefern, wodurch letztere Unbequemlichkeit gehoben und noch der Vortheil gewonnen wird, dass man gleich auf zwei Sterne die beiden Kreise im Voraus stellen und sie also, auch wenn
sie sehr schnell auf einander folgen, doch bequem beobachten kann. Ich theile Ihnen noch meine Beobachtungen der letzten Jupitersopposition mit, wo ich die geraden Aufsteigungen am Mittagsfernrohr
beobachtete, während Hr. Prof. \textsc{Harding} die Zenithdistanzen am \textsc{Repsold}schen Kreise einstellte. Diese
Beobachtungen können jedoch noch nicht als Maassstab für die Genauigkeit, die sich mit dem Mittagsfernrohr erreichen lässt, dienen, weil damals die gebrauchte Uhr einige Unregelmässigkeiten im Gange
zeigte.

\begin{center}
\begin{tabular}{r|ccc|cc|cc}
1819 & \multicolumn{3}{c|}{Mittlere Zeit} & \multicolumn{2}{c|}{} & \multicolumn{2}{c}{} \\
\hline
Juli 30 & $12^{\mathrm{u}}$ & $31^{\mathrm{m}}$ & $46^{\mathrm{s}}4$ & $315^{\circ}\,42'$ & $47''3$ & $-17^{\circ}\,47'$ & $50''8$ \\
Aug.\ 1 & 12 & 22 & 52.5 & $27$ & 13.2 & $52$ & 34.8 \\
3 & 12 & 13 & 58.7 & $11$ & 41.1 & $57$ & 12.1 \\
4 & 12 & 9 & 31.7 & $3$ & 53.5 & $59$ & 35.9 \\
5 & 12 & 5 & 4.5 & $314\,56$ & 3.1 & $18\,1$ & 49.3 \\
13 & 11 & 29 & 29.1 & $313\,53$ & 50.7 & $19$ & 56.1 \\
\end{tabular}
\end{center}

\par\indent Im vorigen Herbst erhielt ich von Hrn. \textsc{Repsold} eine neue höchst vortreffliche Libelle, die unmittelbar auf das Fernrohr gesetzt wird. Die Vergleichung der Beobachtungen, die in entgegengesetzten Lagen des Instruments gemacht werden, gibt den Winkel, welchen die nivellirte Linie mit der Gesichtslinie macht. Seit jener Zeit habe ich die Beobachtungen von etwa 50 Circumpolarsternen zu einem Hauptgeschäft gemacht, und etwa 1200 Beobachtungen erhalten. Die Resultate für die Polhöhe
und die Declinationen dieser Sterne werde ich bei einer andern Gelegenheit bekannt machen, es müssen aber erst noch mehrere andere Prüfungen vorangehen. Theils werde ich erst, mit Hülfe eines besondern Apparats, der jetzt an den Pfeilern angebracht ist, die Theilung auf das sorgfältigste durchprüfen, theils werde ich noch eine andere Reihe von Beobachtungen machen, die ich für unerlässlich
halte, wenn die \emph{absoluten} Zenithdistanzen Zuverlässigkeit erhalten sollen. Vielfache Erfahrungen der
neuern Zeit, an Declinationen von Sternen und bei der Schiefe der Ekliptik haben gezeigt, dass wir \SourceLeaf{432}
über die wahren absoluten Zenithdistanzen noch durchaus in Ungewissheit sind, ein Instrument giebt
sie so, das andere anders. Da wir die Einwirkung der \emph{Schwere} auf die Theile des Instruments
nicht aufheben können, so entstehen nach Maassgabe des Baues des Instruments durch Biegung u. s. w.
\emph{Fehler}, die nach den Zenithdistanzen verschieden sind, und von denen man annehmen darf, dass
sie die Form $a\sin z+b\cos z$ haben. Bei solchen Instrumenten, die sich umwenden lassen, und an denen
jeder Stern in beiden Lagen beobachtet wird, ist der zweite Theil $b\cos z$ unschädlich, er bewirkt
blos einen nach der Zenithdistanz ungleich ausfallenden Collimationsfehler, ohnehin wird aber auch ein
solcher Theil bei einem gut ausgeführten Instrument noch immer merklich sein. Allein anders verhält
es sich mit dem ersten Theil $a\sin z$, alle Zenithdistanzen werden dadurch verfälscht, die kleinen
wenig, die grössern mehr. In dieser Beziehung ist jedes Instrument ein Individuum, und wenn gleich
durch den Künstler dieser Fehler sehr vermindert werden kann, so halte ich es doch für unmöglich,
ihn mit Sicherheit ganz wegzuschaffen, und es bleibt nichts übrig, als ihn durch Beobachtungen zu
bestimmen. Die Vergleichung von den Declinationen südlicher Sterne, welche Herr \textsc{Littrow} aus
\textsc{Bessel}s Beobachtungen abgeleitet hat, mit den \textsc{Pond}schen, zeigt mir, dass die Differenz
der Werthe von $a$ bei \textsc{Pond}s und bei \textsc{Bessel}s Kreise etwa $=3''$ gesetzt werden muss,
um die Resultate am besten in Uebereinstimmung zu bringen, allein wie gross die Werthe einzeln sind,
bleibt unentschieden. Wüsste man \emph{gewiss}, dass der Mittelpunkt der Figur der Sonne einen
grössten Kreis beschreibt, und dass die Refraction für Sonnen- und Sternlicht genau dasselbe ist, so
würde man berechtigt sein, jenen Unterschied ganz oder grösstentheils auf \textsc{Pond}s Kreis
zu schieben; allein da jene Voraussetzungen unerwiesen sind, so würde dieser Schluss eine petitio
principii sein. Ich bemerke nur noch, dass aus Vergleichung meiner Beobachtungen mit den
\textsc{Pond}schen folgt, dass der Werth von $a$ an \textsc{Repsold}s Kreise von dem des
\textsc{Pond}schen um $1''5$ abweicht, und zwar auf der entgegengesetzten Seite von \textsc{Bessel}s
Kreise, so dass dieser vom \textsc{Repsold}schen $4''5$ differirt. Man darf sich darüber gar nicht
wundern, nach mehreren Erfahrungen bewirkt an \textsc{Repsold}s Kreise das Auflegen eines Gewichts
von 1 Loth bei dem Objectiv eine Biegung, wodurch die Zenithdistanzen im Horizont um $1''$ vermindert
werden. Die Balancirungen am Fernrohr, die \textsc{Repsold} und \textsc{Reichenbach} angebracht haben,
machen zwar diese Biegung sehr klein, können aber keine Gewähr leisten, dass sie ganz aufgehoben ist.
Das einzige Mittel, den Werth von $a$ zu erfahren, scheint mir zu sein, dass man eine hinreichend
grosse Anzahl Beobachtungen aus einem Quecksilberhorizont macht. Ich habe bisher nur erst einigemal
Beobachtungen dieser Art gemacht, und zwar am Polarstern; diese geben zwar fast haarscharf im Mittel
$a=0$, allein die Beobachtungen müssten erst viel zahlreicher sein, um mich dabei beruhigen zu
können. Doch kann ich schon glauben, dass der Werth von $a$ an \textsc{Repsold}s Kreise fast $5''$
betragen könnte. Diese Beobachtungen haben übrigens, besonders bei einem grossen Instrument, manche
praktische Schwierigkeit.

\par\indent Von meinen Beobachtungen am \textsc{Repsold}schen Kreise führe ich hier noch die geraden
Aufsteigungen des Nordsterns an, die ich im vorigen Mai bestimmt habe. Der Decimalbruch des Tages ist
von der untern Culmination des Datum an gezählt, 
die Vergleichung ist mit \textsc{Bessel}s \emph{neuen} Tafeln gemacht\footnote{Jedoch mit
Weglassung der beiden kleinen Gleichungen.} (die Unterschiede der Beobachtungen im vorigjährigen
Jahrbuche bezogen sich auf \textsc{Bessel}s ältere Tafeln). Wegen der Columne mit der Ueberschrift
Gewicht, beziehe ich mich auf das vorigjährige Jahrbuch (S. 214 [S. 415 d. B.]).

\Needspace{22\baselineskip}
\SourceLeaf{433}
\begin{center}
\begin{tabular}{r|l|l|l}
1819  &  \shortstack{Scheinb.\\AR.}  &  \shortstack{Verb. Diff. v.\\Bessels Taf.}  &  Gewicht  \\
\hline
 &  $0^{\mathrm u}\,55^{\mathrm m}$  &   &   \\
Mai 1.\,0  &  $51^{\mathrm s}03$  &  $+0^{\mathrm s}07$  &  5.52  \\
2.75  &  51.27  &  $-$0.40  &  3.43  \\
3.25  &  52.39  &  $+$0.51  &  4.44  \\
6.00  &  52.97  &  $-$0.09  &  3.56  \\
7.25  &  54.33  &  $+$0.72  &  3.75  \\
\hline
 1819  &  \shortstack{Scheinb.\\AR.}  &  \shortstack{Verb. Diff. v.\\Bessels Taf.}  &  Gewicht  \\
\hline
  &  $0^{\mathrm u}\,55^{\mathrm m}$  &   &   \\
 Mai 8.25  &  $54^{\mathrm s}83$  &  $+0^{\mathrm s}76$  &  3.43  \\
 9.75  &  54.95  &  $+$0.17  &  3.00  \\
 10.75  &  53.29  &  $-$1.98  &  3.40  \\
 16.50  &  57.21  &  $-$1.06  &  3.00  \\
 21.25  &  60.96  &  $+$0.03  &  2.67  \\
\end{tabular}
\end{center}

Im Mittel wäre hiernach die Correction nur $= -0^{\mathrm s}03$.

\par\indent Seit kurzem ist nun auch der \textsc{Reichenbach}sche Meridiankreis angelangt, welcher auch bald wird aufgestellt werden können. In Beziehung auf das Mittagsfernrohr bemerke ich noch, dass seine optische Kraft ungemein gross ist. Den kleinen Stern, welcher $5'$ vom Pol absteht, sehe ich in der Dämmerung, wo noch keine Fadenbeleuchtung nöthig ist. Das Umlegen geschieht vermittelst eines von dem hiesigen Hrn. Inspector \textsc{Rumpf} sehr schön ausgeführten Apparats in der Luft. Durch ein Seil, welches oben über eine Rolle, von da nach der Seite zu einer zweiten Rolle geht, und dann unten auf die Welle einer Winde mit Sperrung u.\,s.\,w. gewickelt ist, hebt man das Instrument etwa nur 1 Zoll in die Höhe, worauf es freischwebend umgedreht und in der entgegengesetzten Lage wieder eingelassen wird. Die Rolle wird nach gemachtem Gebrauch durch den blossen Zug an einer Schnur seitwärts geschoben, um bei Zenithbeobachtungen nicht zu hindern. Das ganze Geschäft dauert etwa 5--6 Minuten, so dass ich, wenn ich während der Nordstern Culmination umlege, von 7 Fäden nur einen verliere: wenn ich erst den zweiten Stellungskreis erhalten haben werde, wird dies Geschäft leicht so schnell ausgeführt werden können, dass gar kein Faden verloren geht. Beim \textsc{Repsold}schen Kreise kostet das Umlegen immer über eine halbe Stunde. Bei dem \textsc{Reichenbach}schen Kreise wird es auf eine andere Art, vermittelst einer besondern Maschine, geschehen, und der Künstler glaubt, dass 8 Minuten jedesmal zureichend sein werden.
